\documentstyle[%


righttag%


,11pt%


,amssymb%


,russian%               remove in Eng.


,izvru%                 Set izven% in Eng.


]{amsart}





% {,}





\newcommand{\SP}{\operatorname{Sp}}


\newcommand{\Div}{\operatorname{div}}


\newcommand{\tri}{| @!@!@!@!@! | @!@!@!@!@! |}


\newcommand{\norme}[1]{\tri #1 \tri}


\newcommand{\norm}[1]{\| #1 \|}


\newcommand{\abs}[1]{| #1 |}


\newcommand{\tts}[1]{}


\renewcommand{\baselinestretch}{0.96}








%\hoffset=-15mm%                  For Eng.


\setcounter{page}{64}


\god{2005}


\nomer{1~(512)} %                        Remove in Eng.


%\nomer{Vol.\,49, No.\,1}%               Set in Eng.


%\str{\pageref{begin}--\pageref{end}}%  Set in Eng.


\udk{519.635:539.319}





\author{\it А.В.\,Музалевский, С.И.\,Репин}


% NB:   ^^ remove in Eng.


\title{Об оценках погрешности приближенных


решений в задачах линейной теории термоупругости}


\thanks{Работа выполнена при финансовой поддержке гранта Министерства


образования и науки


Российской Федерации (код проекта \No\,Е02-1.0-55).}





\date{}


%\pagestyle{myheadings}%         Set in Eng.


\pagestyle{plain} %              Remove in Eng.


\begin{document}


%\thispagestyle{plain}%          Set in Eng.


\maketitle


\label{begin}














\section{Введение}





Многие реальные проблемы, возникающие в приложениях, приводят к


сложным системам, которые могут включать дифференциальные,


интегральные и алгебраические уравнения. Проблема контроля


точности приближенных решений таких систем весьма сложна,


поскольку каждая из входящих в систему переменных заключает в себе


некоторую погрешность и тем самым индуцирует дополнительные ошибки


в определении всех остальных переменных задачи. В данной


статье рассмотрена одна из типичных задач такого класса, которая


возникает при совместном анализе термических и силовых полей,


образующихся в упругих телах при нагревании. Здесь температурные напряжения,


определяющие дополнительные деформации и напряжения определяются из решения


задачи теплопроводности, присоединенной к уравнениям теории


упругости.





Контроль точности приближенных решений является одной из важнейших задач


вычислительной математики. Эту задачу можно решать как с помощью априорного,


так и апостериорного подхода. Априорные оценки, как правило, требуют 


дополнительной регулярности точного решения и имеют асимптотический 


характер, т.\,е. показывают скорость убывания погрешности при уменьшении 


характерного параметра сетки. Апостериорные оценки возникли несколько 


позднее априорных в связи с необходимостью индикации распределения


погрешности приближенного решения, построенного на конкретной сетке.


Апостериорные индикаторы погрешности широко используются в современных 


вычислительных программах основанных на принципе последовательной 


адаптации сеток. Полное решение задачи апостериорного контроля требует 


не только получения индикатора распределения погрешности, но и построения 


гарантированной верхней границы ошибки аппроксимации в


естественной (энергетической) норме.





Апостериорные оценки ошибки для приближенных решений, полученных методом 


конечных элементов, рассматривались многими авторами. Первым был предложен


так называемый метод ``невязок'' (см., напр., \cite{AiOd}--\cite{Ve}), 


который в дальнейшем получил широкое распространение.


С математической точки зрения этот метод сводится к вычислению нормы 


невязки соответствующего дифференциального уравнения в пространстве 


распределений. Этот метод пригоден только для гал\"еркинских аппроксимаций. 


Позднее для индикации распределения ошибок конечноэлементных решений


был предложен метод осреднения градиента (см., напр.,


\cite{Wa}, \cite{ZiZh1}). Этот метод также применим только для 


гал\"еркинских аппроксимаций и требует повышенной регулярности точного


решения задачи.





Мажоранты энергетической нормы отклонения от точного решения, которые пригодны 


для любых функций из энергетического класса, были получены в


\cite{Re1}--\cite{Re3}. Эти мажоранты получены из общих методов 


функционального анализа и теории двойственности для вариационного исчисления,


поэтому их можно называть апостериорными оценками функционального типа. 


Для задач линейной


теории упругости они ранее изучались в работах \cite{MuRe}, \cite{Re3}.





В данной работе этот подход используется для получения мажоранты 


функционального типа для задач линейной теории термоупругости и проверяется


его эффективность на серии тестовых примеров.








\section{Вариационные постановки}





Пусть $\Omega$ --- ограниченная связная область в ${\Bbb R}^n$ c


липшицевой границей $\partial \Omega$, причем $\partial \Omega$


состоит из двух непересекающихся частей $\partial_1 \Omega$ и


$\partial_2 \Omega$.


Классическая постановка краевой задачи теории термоупругости


состоит в определении тензор-функции $\sigma$ и вектор-функции


$u$, удовлетворяющих следующей системе уравнений:


\begin{gather}


\Div \sigma + f + f_T =0 \quad \text{в} \quad \Omega,


\label{el1}\\


%


\sigma = {\Bbb L} \varepsilon (u) \quad \text{в} \quad \Omega,


\label{oe1}\\


%


\varepsilon (u) = \frac{1}{2} (\nabla u + (\nabla u)^T), \\


%


u=u_0  \quad \text{на} \quad \partial_1\Omega, \\


%


\sigma\nu =F+F_T \quad \text{на} \quad \partial_2\Omega,


\label{oe5}


\end{gather}


где $F$, $u_0$ --- заданные функции, определяющие краевые условия


задачи упругости, $\varepsilon(u)$ обозначает тензор малых


деформаций, $f$ --- заданная функция, определяющая объемные силы,


$\nu$ --- внешняя нормаль к границе области,


${\Bbb L}=\{{\Bbb L}_{ijkl}\}$ --- тензор упругих


модулей, который должен удовлетворять условию


\begin{equation*}


c_1\abs{\varkappa}^2 \le {\Bbb L}\varkappa : \varkappa \le c_2\abs{\varkappa}^2


\quad \forall \varkappa\in {\Bbb M}^{n\times n}_s


%% \label{ineq:bfl}


\end{equation*}


и условиям симметричности


$$


{\Bbb L}_{ijkl}={\Bbb L}_{jikl}={\Bbb L}_{klij}\quad


\forall i,j,k,l=1,\dotsc,n,


$$


а $f_T$, $F_T$ ---


функции, связанные с тензором температурных деформаций


$$


\varepsilon_T =


\begin{pmatrix}


\alpha_{11} & 0  & 0 \\


0 & \alpha_{22} & 0 \\


0 & 0  & \alpha_{33}


\end{pmatrix} (T-T_0)


$$


соотношениями (см., напр., \cite{TiGu}, с.\,458)


$f_T=-\Div({\Bbb L}\varepsilon_T)$,


$F_T={\Bbb L}\varepsilon_T$.


Здесь $\alpha_{11}$, $\alpha_{22}$, $\alpha_{33}$ --- 


коэффициенты температурного


расширения по соответствующим координатным осям, $T_0$ --- отсчетная


температура (т.\,е. температура, при которой отсутствуют температурные


напряжения), а скалярная функция $T$ является решением задачи теплопроводности


\begin{equation}\label{th1}


-\Delta T = q


\end{equation}


с граничными условиями


\begin{equation}


\label{th2}


T=\theta_0\ \ \text{на}\ \ \partial_3\Omega, \quad


\frac{\partial T}{\partial \nu} = Q \ \ \text{на} \ \ \partial_4\Omega.


\end{equation}


Здесь $q$ --- функция, определяющая тепловые источники в области $\Omega$,


$\partial_3 \Omega$ и $\partial_4 \Omega$ --- две


непересекающиеся части границы $\partial \Omega$, которые, вообще


говоря, могут не совпадать с $\partial_1 \Omega$ и


$\partial_2\Omega$, $\theta_0$, $Q$ --- заданные функции, определяющие


соответственно заданную температуру на границе $\partial_3 \Omega$


и тепловой поток через $\partial_4 \Omega$.





Нетрудно видеть, что погрешности аппроксимации, возникающие в задаче


(\ref{th1}), (\ref{th2}) вносят дополнительные погрешности в численное 


решение основной задачи (\ref{el1})--(\ref{oe5}).


Для того чтобы учесть суммарный эффект такого наложения ошибок, мы используем


подход, отличный от того, что обычно используется для гал\"еркинских 


аппроксимаций (см., напр., \cite{AiOd}--\cite{Ve}). Наш подход основан 


на использовании апостериорных оценок функционального типа


(см. \cite{Re3}, \cite{Re4}).





Введем следующие обозначения:


\begin{gather*}


V_0=\{ u\in W_2^1(\Omega,{\Bbb R}^n) \mid u=0


\ \text{на} \ \partial_1\Omega\},\\


%


U_0=\{ u\in W_2^1(\Omega) \mid u=0 \ \text{на} \ \partial_3\Omega\},\\


%


\norme{\varepsilon(u)}^2 =\int_{\Omega}


{\Bbb L} \varepsilon (u) : \varepsilon(u) dx,


\quad (\varepsilon, \tau) =\int_{\Omega}


\varepsilon  : \tau\, dx,\\


%


V_0+u_0=\{ u\in W_2^1(\Omega,{\Bbb R}^n) \mid u=u_0 \ \text{на}\


\partial_1\Omega\},\\


%


Y^*=L_2(\Omega,{\Bbb M}^{n\times n}_s).


\end{gather*}


Данная задача может быть сформулирована в виде двух вариационных проблем,


которые будем называть задачами $\cal P$ и $\cal Q$.


\medskip





{\it Задача $\cal P$.} Найти элемент $u \in V_0+u_0$ такой, что


$$


J(u)= \inf_{v\in V_0+u_0} J(v),


$$


где


$$


J(v)=\frac{1}{2} \norme{\varepsilon(v)}^2-\int_{\Omega} (f+f_T)\cdot v \,dx


-\int_{\partial_2\Omega} (F+F_T) \cdot v \, d\Gamma.


$$





{\it Задача $\cal Q$.} Найти элемент $T \in U_0+\theta_0$ такой, что


$$


I(T)= \inf_{\theta\in U_0+\theta_0} I(\theta),


$$


где


$$


I(\theta)=\frac{1}{2} \norm{\nabla\theta}^2-\int_{\Omega} q \theta \,dx


-\int_{\partial_4\Omega} Q  \theta \, d\Gamma.


$$


Пусть $v\in V_0+u_0$ и $\theta \in U_0+\theta_0$ являются приближенными 


решениями задачи термоупругости. Нашей задачей является получение 


апостериорной оценки для $u-v$.





Методика получения апостериорных оценок в энергетических нормах для широкого


класса эллиптических краевых задач предложена в \cite{Re3}, \cite{Re4}. 


Для задач линейной упругости эти оценки изучались в \cite{MuRe}, где 


было продемонстрировано


их эффективное использование в практических вычислениях.





Если положить $F_T\equiv 0$ и $f_T\equiv 0$, то оценка


для получившейся задачи теории упругости может быть


записана в следующем виде:


\begin{multline}


\frac{1}{2}\norme{\varepsilon(v-u)}^2 \le


{\cal M}^{(1)}_\oplus(v,\beta,\tau) =


\frac{1+\beta}{2}


\big(\varepsilon(v) -{\Bbb L}^{-1}\tau, {\Bbb


L}\varepsilon(v) - \tau\big) +\\


%


+\frac{1}{2}\bigg(1+\frac{1}{\beta} \bigg) C_{1\Omega}^2


(\norm{\Div \tau + f}^2_{\Omega} +


\norm{F-\tau\nu}^2_{\partial_2\Omega}).


\label{mpluse}


\end{multline}


Здесь $\tau \in Y^*$ и $\beta>0$ --- произвольные параметры, а


$C_{1\Omega}$ --- константа, определяемая неравенством


\begin{equation}


\label{ineq:c1}


\norm{w}^2_{\Omega}  + \norm{w}^2_{\partial_2\Omega}


\le C_{1\Omega}^2 \norme{\varepsilon(w)}^2 \quad \forall w\in V_0.


\end{equation}





В свою очередь оценка для задачи теплопроводности выглядит


следующим образом:


\begin{multline}


\frac{1}{2}\norm{\nabla(\theta-T)}^2_{\Omega} \le


{\cal M}^{(2)}_\oplus(\theta,\beta,y) =


\frac{1+\beta}{2} (\nabla \theta - y,


\nabla \theta - y ) +\\


%


+\frac{1}{2} \bigg(1+\frac{1}{\beta} \bigg)


C_{2\Omega}^2 \bigg(\norm{\Div y + q}^2_{\Omega} +


\norm{\frac{\partial y}{\partial \nu} - Q}^2_{\partial_4\Omega} \bigg),


\label{mplust}


\end{multline}


где $y\in L_2(\Omega,{\Bbb R}^n)$ и $\beta>0$ --- произвольные


параметры, а $C_{2\Omega}$ --- константа, определяемая аналогичным


(\ref{ineq:c1}) неравенством


\begin{equation*}


%% \label{ineq:c2}


\norm{w}^2_{\Omega}  + \norm{w}^2_{\partial_4\Omega}


\le C_{2\Omega}^2 \norm{\nabla w}^2_{\Omega} \quad \forall w\in U_0.


\end{equation*}





Заметим, что если $\partial_4\Omega=\emptyset$, то константа $C_{2\Omega}$


является константой в неравенстве Фридрихса. В этом случае ее можно


оценить аналитически как


$$


C_{2\Omega}\le \frac{l}{\sqrt{2}\pi},


$$


где $l$ --- сторона квадрата, охватывающего область $\Omega$.








В \cite{Re3} было показано, что, во-первых, последовательность этих оценок,


полученная путем решения только конечномерных задач, будет


сходиться к точной оценке с увеличением размерности


соответствующих конечномерных пространств, и, во-вторых,


выбором свободных переменных $\beta$ и $\tau$ (или


соответственно $\beta$ и $y$) можно получить


равенство оценки точной величине отклонения.





В следующем параграфе данные оценки используются для построения апостериорной


оценки в термоупругой задаче.








\section{Оценки отклонения от точного решения в энергетической норме}





Применив оценку (\ref{mpluse}), приходим к неравенству


\begin{multline}


\frac{1}{2}\norme{\varepsilon(v-u)}^2 \le \frac{1+\beta_1}{2}


\big(\varepsilon(v) -{\Bbb L}^{-1}\tau, {\Bbb


L}\varepsilon(v) - \tau\big) +\\


%


+\frac{1}{2}\bigg(1+\frac{1}{\beta_1} \bigg) C_{1\Omega}^2


(\norm{\Div \tau + f + f_T}^2_{\Omega} +


\norm{F+F_T-\tau\nu}^2_{\partial_2\Omega}),


\label{mplus1}


\end{multline}


которое справедливо для любых $\tau \in Y^*$ и $\beta_1>0$.


Заметим, что в правую часть (\ref{mplus1}) входят неизвестные


величины $f_T$, $F_T$, которые зависят от точного решения задачи


теплопроводности. Слагаемые, включающие в себя неизвестные величины,


можно оценить через неравенства


\begin{gather*}


\norm{\Div \tau + f + f_T}^2_{\Omega}\le 


(1+\beta_2) \norm{\Div \tau + f + f_{\theta}}^2_{\Omega}


+\bigg(1+\frac{1}{\beta_2}\bigg)


\norm{f_T-f_\theta}^2_\Omega,\\


%


\norm{F+F_T-\tau\nu}^2_{\partial_2\Omega} \le 


(1+\beta_3) \norm{F+F_\theta-\tau\nu}^2_{\partial_2\Omega}


+\bigg(1+\frac{1}{\beta_3}\bigg)


\norm{F_T-F_\theta}^2_{\partial_2\Omega},


\end{gather*}


верные для любых положительных $\beta_2$ и $\beta_3$.





Теперь необходимо оценить величины $\norm{f_T-f_\theta}^2_\Omega$ и


$\norm{F_T-F_\theta}^2_{\partial_2\Omega}$, где


$$


f_\theta = -\Div \left(\Bbb L


\begin{pmatrix}


\alpha_{11} & 0  & 0 \\


0 & \alpha_{22} & 0 \\


0 & 0  & \alpha_{33}


\end{pmatrix} (\theta-T_0)\right),\quad


F_\theta = \Bbb L


\begin{pmatrix}


\alpha_{11} & 0  & 0 \\


0 & \alpha_{22} & 0 \\


0 & 0  & \alpha_{33}


\end{pmatrix} (\theta-T_0).


$$


Заметим, что ошибки в силовых добавках удобно будет оценить


через $\norm{\nabla(\theta-T)}_\Omega$, поскольку оценка этой


ошибки уже известна как ошибка задачи теплопроводности (\ref{mplust}).


Рассмотрим неравенства


$$


\norm{f_T-f_\theta}_\Omega \le


C_{3\Omega} \norm{\nabla(\theta-T)}_\Omega, \quad


\norm{F_T-F_\theta}_{\partial_2\Omega} \le


C_{4\Omega} \norm{\nabla(\theta-T)}_\Omega.


$$


Здесь $C_{3\Omega}$ --- константа, зависящая от термоупругих свойств


среды, которая может быть получена просто из выражения для $f_\theta$.


Действительно, обозначив


$$


\xi=


\begin{pmatrix}


\alpha_{11} & 0  & 0 \\


0 & \alpha_{22} & 0 \\


0 & 0  & \alpha_{33}


\end{pmatrix},


$$


получим


$f_\theta = -\Div (\Bbb L \xi)(\theta-T_0)- (\Bbb L \xi) \nabla\theta$.


Откуда


$$


\norm{f_T-f_\theta}_\Omega \le\sup\limits_\Omega


\abs{\Div(\Bbb L \xi)}\,\norm{\theta-T}+


\sup\limits_\Omega\abs{\Bbb L \xi}\,\norm{\nabla(\theta-T)}


\ \ \text{и} \ \


C_{3\Omega}=C_\Omega\sup\limits_\Omega


\abs{\Div(\Bbb L \xi)}+\sup\limits_\Omega\abs{\Bbb L \xi},


$$


где $C_\Omega$ --- константа в неравенстве Фридрихса.


Очевидно, в случае изотропной модели среды, константа $C_{3\Omega}$


не зависит от $\Omega$. С другой стороны,


$$


C_{4\Omega}=\widehat C_{4\Omega} \sup_\Omega\abs{\Bbb L \xi},


$$


где $\widehat C_{4\Omega}$ является нормой оператора следа на


$\partial_2\Omega$,


$$


\norm{\varphi}_{\partial_2\Omega} \le


\widehat C_{4\Omega} \norm{\nabla\varphi}_\Omega \quad


\forall \varphi\in W_1^2(\Omega).


$$





Собирая все оценки вместе, получим неравенство


\begin{multline*}


\frac{1}{2}\norme{\varepsilon(v-u)}^2 \le


{\cal M}^{(3)}_\oplus(v,\theta,\tau,y,\beta_1,\beta_2,\beta_3,\beta_4)=\\


%


=\frac{1}{2}(1+\beta_1) \big(\varepsilon(v) -{\Bbb


L}^{-1}\tau, {\Bbb L}\varepsilon(v) - \tau\big)+\\


%


+\frac{1}{2}\left(1+\frac{1}{\beta_1}\right) C_{1\Omega}^2 \Bigg[


(1+\beta_2)\norm{\Div \tau + f + f_\theta}^2_{\Omega}+


(1+\beta_3)\norm{F+F_\theta-\tau\nu}^2_{\partial_2\Omega} +\\


%


+\bigg(C_{3\Omega}^2\bigg(1+\frac{1}{\beta_2}\bigg)


+C_{4\Omega}^2\bigg(1+\frac{1}{\beta_3}\bigg)\bigg) \bigg[


(1+\beta_4) \norm{\nabla \theta - y}^2 +\\


%


+\bigg(1+\frac{1}{\beta_4} \bigg) C_{2\Omega}^2


\bigg(\norm{\Div y + q}^2_{\Omega} + \bigg\|\frac{\partial


y}{\partial \nu} - Q\bigg\|^2_{\partial_4\Omega} \bigg)\bigg]\Bigg],


\end{multline*}


которое верно для любых $\tau\in Y^*$, $y\in L_2(\Omega,{\Bbb R}^n)$


и любых положительных $\beta_1$, $\beta_2$, $\beta_3$, $\beta_4$.





Если задача термоупругости ставится при


$\partial \Omega_2 =\emptyset$, $\partial\Omega_4=\emptyset$, то


выражение для мажоранты упрощается, т.\,к. пропадают слагаемые, отвечающие за


ошибки на границе. В этом случае соответствующая оценка имеет вид


\begin{gather*}


\frac{1}{2}\norme{\varepsilon(v-u)}^2 \le


{\cal M}^{(4)}_\oplus(v,\theta,\tau,y,\beta_1,\beta_2,\beta_3)=


\frac{1}{2}(1+\beta_1)


\big(\varepsilon(v) -{\Bbb L}^{-1}\tau,


{\Bbb L}\varepsilon(v) - \tau\big)+\\


%


+\frac{1}{2}\bigg(1+\frac{1}{\beta_1}\bigg) C_{1\Omega}^2 \Bigg[


(1+\beta_2)\norm{\Div \tau + f + f_\theta}^2_{\Omega}+ \\


%


+C_{3\Omega}^2 \bigg(1+\frac{1}{\beta_2}\bigg)


\bigg[ (1+\beta_3) \norm{\nabla \theta - y}^2


+ \bigg(1+\frac{1}{\beta_3} \bigg)


C_{2\Omega}^2 (\norm{\Div y + q}^2_{\Omega}


)\bigg]\Bigg].


\end{gather*}








\section{Оценки для изотропной модели}





Для трехмерной изотропной модели линейной теории упругости


тензор упругих постоянных ${\Bbb L}$ определяется двумя постоянными


$K_0$ и $\mu$, которые называются соответственно модулями


объемного сжатия и сдвига. В этом случае


определяющие соотношения (\ref{oe1}) имеют вид


\begin{equation*}


\varepsilon (u)=\frac{1}{9 K_0}\SP\sigma {\bold 1}


+ \frac{1}{2\mu} \sigma^D


\end{equation*}


или


$$


\varepsilon (u)=\frac{1-2\nu}{3E} \SP\sigma {\bold 1}


+ \frac{1+\nu}{E} \sigma^D,


$$


где $\bold 1$ --- единичный тензор, $\SP\sigma := {\bold 1} : \sigma$


--- след тензора, $\sigma^D=\sigma-\frac{1}{3} {\bold 1} \SP\sigma$


--- девиатор тензора $\sigma$,


а $E$, $\nu$ --- часто используемые в прикладных задачах параметры,


определяющие упругие свойства среды (модуль Юнга и


коэффициент Пуассона). Они связаны с $K_0$ и $\mu$ соотношениями


$$


K_0=\frac{E}{3(1-2\nu)}, \quad \mu = \frac{E}{2(1+\nu)}.


$$


Для изотропного тела $\alpha_{ii}=\alpha$, $i=1,2,3$,


$$


f_T=-\frac{\alpha E}{1-2\nu} \nabla T, \quad


F_T=\frac{\alpha E}{1-2\nu} (T-T_0).


$$


Введем обозначение $\beta=\frac{\alpha E}{1-2\nu}$,


тогда, принимая во внимание определяющие соотношения


$$


{\Bbb L}\varepsilon=K_0 \SP\varepsilon {\bold 1}+ 2\mu \varepsilon^D,


$$


получим $C_{3\Omega}=\sup\limits_\Omega\abs{{\Bbb L}\xi}=\beta$,


$C_{4\Omega}=\widehat C_{4\Omega} \beta$.


Oценку (\ref{mplus1}) можно записать в виде


\begin{gather*}


\frac{1}{2}\norme{\varepsilon(v-u)}^2 \le


{\cal M}_\oplus(v,\theta,\tau,y,\beta_1,\beta_2,\beta_3,\beta_4)= 


\frac{1}{2}(1+\beta_1)


\big(\varepsilon(v) -{\Bbb L}^{-1}\tau,


{\Bbb L}\varepsilon(v) - \tau\big)+ \\


%


+\frac{1}{2}\bigg(1+\frac{1}{\beta_1}\bigg) C_{1\Omega}^2 \Bigg[


(1+\beta_2)\norm{\Div \tau + f - \beta\nabla\theta}^2_{\Omega}+ 


(1+\beta_3)\norm{F+\beta(\theta-T_0)-\tau\nu}^2_{\partial_2\Omega} + \\


%


+\beta^2\bigg(\bigg(1\!+\!\frac{1}{\beta_2}\bigg)


+\widehat C_{4\Omega}^2 \bigg(1\!+\!\frac{1}{\beta_3}\bigg)\bigg)


\bigg[ (1\!+\!\beta_4) \norm{\nabla t\! -\! y}^2\!+\!


%%%%%%%% pay attn to \! and extra line


\bigg(1\!+\!\frac{1}{\beta_4} \bigg)


C_{2\Omega}^2 \bigg(\norm{\Div y\!+\!q}^2_{\Omega}\!+\!


\bigg\|\frac{\partial y}{\partial \nu}\!-\!Q\bigg\|^2_{\partial_4\Omega}


\bigg)\bigg]\Bigg].


\end{gather*}


Здесь $\tau\in Y^*$ --- произвольная тензор-функция,


$y\in L_2(\Omega,{\Bbb R}^n)$ --- произвольная вектор-функция,


$\beta_1$, $\beta_2$, $\beta_3$, $\beta_4$ --- положительные константы.


Заметим, что эта оценка обращается в нуль только тогда, когда


$v=u$ и $\theta=T$, т.\,е. для точного решения задачи теплопроводности.


Также обратим внимание, что оценку легко разбить на слагаемые, каждое


из которых отвечает за ошибки в разных уравнениях классической постановки,


т.\,е. первое слагаемое отвечает за ошибку в определяющих соотношениях,


второе и третье --- за ошибки в силовых условиях, четвертое, пятое и шестое


--- за ошибки в задаче теплопроводности.





Для первой краевой задачи термоупругости, которая получается при


$\partial \Omega_2 =\emptyset$ и $\partial\Omega_4=\emptyset$,


выражение для мажоранты упрощается. В этом случае


\begin{gather*}


\frac{1}{2}\norme{\varepsilon(v-u)}^2 \le


{\cal M}_\oplus(v,\theta,\tau,y,\beta_1,\beta_2,\beta_3)=


\frac{1}{2}(1+\beta_1)


\big(\varepsilon(v) -{\Bbb L}^{-1}\tau,


{\Bbb L}\varepsilon(v) - \tau\big)+ \\


%


+\frac{1}{2}\bigg(1+\frac{1}{\beta_1}\bigg) C_{1\Omega}^2 \Bigg[


(1+\beta_2)\norm{\Div \tau + f - \beta \nabla \theta}^2_{\Omega}+\\


%


+\beta^2 \bigg(1+\frac{1}{\beta_2}\bigg)


\bigg[ (1+\beta_3) \norm{\nabla \theta - y}^2


+ \bigg(1+\frac{1}{\beta_3} \bigg)


C_{2\Omega}^2 (\norm{\Div y + q}^2_{\Omega}


)\bigg]\Bigg].


\end{gather*}








\section{Численные эксперименты}





Полученные выше оценки были использованы для анализа качества приближенных 


решений ряда задач термоупругости. Ниже мы приводим


некоторые из полученных результатов, в которых расчеты


производились для модели плоского напряженного состояния


с постоянными $E=4$, $\nu=0{,}3$ и $\alpha=1$.





Качество получаемых оценок принято


оценивать при помощи так называемого индекса эффективности


$I_{\text{eff}}$, который определяется следующим образом:


$$


I_{\text{eff}}^2=\frac{2 {\cal M}_\oplus(v,\theta,\tau,y,\beta_1,\beta_2,


\beta_3)}{\norme{\varepsilon(v-u)}^2}.


$$


Ясно, что чем ближе $I_{\text{eff}}$ к единице, тем выше качество


полученной оценки.








Для последующей адаптации сетки важно


знать не только интегральное значение мажоранты, но и величину


вклада в него на каждом элементе разбиения. Так как мажоранта


${\cal M}_\oplus (v,\beta,\tau)$ представляет собой


интеграл, то возникает вопрос о том, насколько хорошо соответствующие


подинтегральные выражения воспроизводят локальное распределение


ошибок по области $\Omega$. Этот вопрос также исследовался в


приведенных ниже экспериментах.








\begin{exam}


Рассматривается задача в квадрате $(0,1)\times (0,1)$.





В табл.~1 приведены соответствующие результаты. Здесь


$N_{\text{elm}}$ обозначает число


элементов в триангуляции,


$M_\oplus$ --- значение мажоранты, полученной в результате численного


эксперимента, $M^{(1)}$, $M^{(2)}$, $M^{(3)}$ --- вклады в мажоранту


ошибки в определяющих соотношениях, ошибки в уравнениях равновесия и ошибки


задачи теплопроводности соответственно,


$I_{\text{eff}}$ --- индекс эффективности.


На рис.\,1 приведены трехмерные гистограммы точной ошибки и


мажоранты ($N_{\text{elm}}=312$).


\smallskip





\begin{center}


Таблица 1. Результаты для примера 1


\smallskip





\begin{tabular}{|c|c|c|c|c|c|c|}


\hline


$N_{\text{elm}}$ & $M^{(1)}$ & $M^{(2)}$


& $M^{(3)}$ & $M_\oplus$ & $\frac{1}{2}\norme{\varepsilon(v-u)}^2$


& $I_{\text{eff}}$ \\ \hline


84 & 0,003110 & 0,001457 & 0,001447 & 0,006015 & 0,001274 & 2,17 \\ \hline


312 & 0,000911 & 0,000542 & 0,000443 & 0,001896 & 0,000369 & 2,26 \\ \hline


1372 & 0,000187 & 0,000110 & 0,000087 & 0,000384 & 0,000084 & 2,13 \\ \hline


5458 & 0,000047 & 0,000029 & 0,000021 & 0,000097 & 0,000021 & 2,14 \\ \hline


\end{tabular}


\end{center}





\begin{center}


\unitlength=1mm


\begin{picture}(170,166) %% width, heigth, added 5 mm for signature


\put(27,166){\special {em:graph mz-pic1.bmp}} %% left X, upper Y, -, /2, +toX


\put(80,0){\mbox Рис.\,1}


\end{picture}


\end{center}


\end{exam}





\begin{exam}


Пусть $\Omega$ представляет собой область


$(-1,1)\times (-1,1)$, в которой сделаны боковые


треугольные вырезы. Численное моделирование производилось для


одной четвертой этой области.





Соответствующие результаты приведены в табл.~2. На рис.\;2


приведены трехмерные гистограммы точной ошибки и


мажоранты ($N_{\text{elm}}=286$) для задачи термоупругости.


\smallskip





\begin{center}


Таблица 2. Результаты для примера 2


\smallskip





\begin{tabular}{|c|c|c|c|}


\hline


$N_{\text{elm}}$ & $M_\oplus$ &


$\frac{1}{2}\norme{\varepsilon(v-u)}^2$ & $I_{\text{eff}}$ \\ \hline


82  & 0,00348 & 0,00154 & 1,50 \\ \hline


286  & 0,00130 & 0,00042 & 1,76 \\ \hline


1131 & 0,00050 & 0,00012 & 2,03 \\ \hline


\end{tabular}


\end{center}





\begin{center}


\unitlength=1mm


\begin{picture}(170,170) %% width, heigth, added 5 mm for signature


\put(29,170){\special {em:graph mz-pic2.bmp}} %% left X, upper Y, -, /2, +toX


\put(80,0){\mbox Рис.\,2}


\end{picture}


\end{center}





%\begin{center}


%\unitlength=1mm


%\begin{picture}(170,116) %% width, heigth, added 5 mm for signature


% \put(2,116){\special {em:graph fig4.bmp}} %% left X, upper Y, -, /2, +toX


%\put(80,0){\mbox Рис.\,2}


%\end{picture}


%\end{center}





Из этих результатов видно, что функционал $M_\oplus$ дает


хорошие оценки величины погрешности,


а интегранд мажоранты достаточно точно воспроизводит распределение


ошибки по области.


\end{exam}








\begin{thebibliography}{393}





\bibitem{AiOd} %1


Ainsworth M., Oden J.T. {\em A posteriori error estimation in finite element 


analysis}. -- John Wiley \& Sons, Inc., 2000. -- 264 p.





\bibitem{BaSt} %2


Babu\v ska I., Strouboulis T. {\em The finite element method and its


reliability}. --


Oxford University Press Inc., New York, 2001. -- 736 p.





\bibitem{Ve} %3


Verf\"urth R. {\em A review of a posteriori error estimation and adaptive


mesh-refinement techniques}. -- Wiley--Teubner, 1996. -- 127 p. %% (+vi)





\bibitem{Wa} %4


Wahlbin L.B. {\em Superconvergence in Galerkin finite element methods} //


Lect. Notes Math., Springer-Verlag. -- 1995. -- V.\,1605.





\bibitem{ZiZh1} %5


Zienkeiewicz O.C., Zhu J.Z. {\em A simple error estimator and adaptive 


procedure for practical engineering analysis} //


Internat. J. Numer. Methods Engrg. -- 1987. -- Т.\,24. -- С.\,337--357.





\bibitem{Re1} %6


Repin S. {\em A posteriori error estimation for nonlinear variational


problems by duality theory} // Zapiski Nauchnych Seminarov


V.A.\,Steklov Math. Inst. in


St.-Petersburg. -- 1997. -- Т.\,243. -- С.\,201--214.





\bibitem{Re2} %7


Repin S. {\em A posteriori error estimation for variational problems


with uniformly convex functionals} // Math. Comput. -- 2000. --


Т.\,69. -- С.\,481--500.





\bibitem{Re3} %8


Репин С.И. {\em Двусторонние оценки отклонения от точного


решения для равномерно эллиптических уравнений} //


Тр. Санкт-Петерб. матем. о-ва. -- 2001. -- Т.\,9. -- С.\,148--179.





\bibitem{MuRe} %9


Muzalevsky A.V., Repin S.I. {\em On two-sided error estimates


for approximate solutions of problems in the linear theory


of elasticity} // Russ. J. Numer. Anal. Math. Modelling. --


2003. -- Т.\,18. -- С.\,65--85.





\bibitem{TiGu} %10


Тимошенко С.П., Гудьер Дж. {\em Теория упругости}. -- М.: Наука, 1979.


-- 576 с.





\bibitem{Re4} %11


Repin S. {\em A unified approach to a posteriori error estimation based 


on duality error majorants} // Math. Comput. Simul. -- 1999. --


Т.\,50. -- С.\,305--321.








\end{thebibliography}





\vskip 2cm





\post{\it Санкт-Петербургский государственный}{\it Поступила}


\post{\it технический университет}{17.09.2004}


\smallskip


\post{\it Санкт-Петербургское отделение}{}


\post{\it математического института им. В.А.\,Стеклова}{}


\post{\it Российской Академии наук}{}





\label{end}


\end{document}








195251, С.-Петербург, Политехническая 29}


%\email{ alxmuz@mail.ru, muz@amd.stu.neva.ru}


% phone  (812) 552-85-92





191011, С.-Петербург, Фонтанка, 27}


%\email{repin@pdmi.ras.ru}


% phone  (812) 535-34-47











